\chapter*{СПИСОК СОКРАЩЕНИЙ И УСЛОВНЫХ ОБОЗНАЧЕНИЙ}
\addcontentsline{toc}{chapter}{СПИСОК СОКРАЩЕНИЙ И УСЛОВНЫХ ОБОЗНАЧЕНИЙ}
\setlength{\parindent}{0pt}
% Перечень располагается столбцом, слева без абзацного отступа в
% алфавитном порядке,
% справа через тире — расшифровка, без знаков препинания в конце
% строки (п. 4.8.4)
BPE - Byte Pair Encoding (кодирование байтовых пар)\\[0.5em]
BPN - Bits Per Nucleotide (бит на нуклеотид)\\[0.5em]
CLM - Causal Language Modeling (каузальное языковое моделирование)\\[0.5em]
CR - Compression Ratio (коэффициент сжатия)\\[0.5em]
dDP - Differentiable Dynamic Programming (дифференцируемое
динамическое программирование)\\[0.5em]
GBST - Gradient-Based Subword Tokenization\\[0.5em]
GPU - Graphics Processing Unit (графический процессор)\\[0.5em]
LM - Language Model (языковая модель)\\[0.5em]
MLM - Masked Language Modeling (маскированное языковое моделирование)\\[0.5em]
MLP - Multilayer Perceptron (многослойный перцептрон)\\[0.5em]
NTP - Next Nucleotide Prediction (предсказание следующего нуклеотида)\\[0.5em]
PE - Positional Encoding (позиционное кодирование)\\[0.5em]
PPL - Perplexity (перплексия)\\[0.5em]
PWM - Position Weight Matrix (позиционная весовая матрица)\\[0.5em]
RC - Reverse Complement (обратный комплемент)\\[0.5em]
RoPE - Rotary Positional Embedding\\[0.5em]
SSM - State Space Model (модель пространства состояний)\\[0.5em]
STE - Straight-Through Estimator\\[0.5em]
TF - Transcription Factor (транскрипционный фактор)\\[0.5em]
TFBS - Transcription Factor Binding Site (сайт связывания
транскрипционного фактора)\\[0.5em]
TR - Token Ratio (доля токенов к длине исходной последовательности)\\[0.5em]
ВКР - выпускная квалификационная работа\\[0.5em]
ДНК - дезоксирибонуклеиновая кислота
